\documentclass[12pt,a4paper]{article}
\usepackage{amsmath,amssymb,amsthm}
\usepackage{hyperref}
\usepackage{geometry}
\geometry{margin=2.5cm}
\usepackage{enumitem}
\usepackage{booktabs}

\newtheorem{theorem}{Theorem}[section]
\newtheorem{lemma}[theorem]{Lemma}
\newtheorem{corollary}[theorem]{Corollary}
\newtheorem{proposition}[theorem]{Proposition}
\theoremstyle{definition}
\newtheorem{definition}[theorem]{Definition}
\theoremstyle{remark}
\newtheorem{remark}[theorem]{Remark}

\title{Biological Homochirality from Recognition Science:\\
$\varphi$-Cascade Symmetry Breaking and the Origin of Life's Handedness}

\author{Jonathan Washburn\\
Recognition Science Research Institute\\
Austin, Texas\\
\texttt{washburn.jonathan@gmail.com}}

\date{March 2026}

\begin{document}
\maketitle

\begin{abstract}
We derive biological homochirality---the exclusive use of L-amino acids
and D-sugars in all known life---from the Recognition Science (RS)
framework with zero free parameters.  The cost functional
$J(x) = \tfrac{1}{2}(x + x^{-1}) - 1$ is symmetric
($J(x) = J(x^{-1})$), so neither chirality is intrinsically
preferred.  However, the $\varphi$-forcing equation ($x^2 = x + 1$)
breaks this symmetry dynamically: the golden ratio satisfies
$\varphi - \varphi^{-1} = 1$ exactly, creating a cascade amplification
in which the $\varphi^n$ branch grows without bound while
$\varphi^{-n} \to 0$.  A tiny initial bias (e.g., from parity
violation at $\sim 10^{-17}$~eV) is amplified exponentially by this
$\varphi$-cascade.  The correlation between L-amino acids and D-sugars
is forced by complementary recognition: $\varphi \times \varphi^{-1}
= 1$ ensures proteins and nucleic acids have equal $J$-cost and can
bind each other.  \end{abstract}

%% ============================================================
\section{Recognition Science Framework}
\label{sec:framework}
%% ============================================================

The homochirality derivation rests on two properties of the cost
functional: its mirror symmetry $J(x) = J(x^{-1})$, and the dynamical
asymmetry of the $\varphi$-cascade that breaks this symmetry.

\begin{definition}[RS Cost Functional]
For $x > 0$: $J(x) = \tfrac{1}{2}(x + x^{-1}) - 1$.
\end{definition}

\noindent\textbf{Axiom A1 (Normalization):} $J(1) = 0$.\\
\textbf{Axiom A2 (Recognition Composition Law, RCL):}
$J(xy)+J(x/y)=2J(x)J(y)+2J(x)+2J(y)$ for all $x,y>0$.\\
\textbf{Axiom A3 (Calibration):} $J''_{\log}(0) = 1$, where
$J_{\log}(t) := J(e^t)$.

\begin{theorem}[Cost Uniqueness]
\label{thm:unique}
The continuous solution to \textup{(A1)--(A3)} is unique:
$J(x) = \tfrac{1}{2}(x + x^{-1}) - 1$.
\end{theorem}

\begin{proof}[Proof sketch]
Set $H(t) = J_{\log}(t)+1$.  Axiom A2 transforms into the d'Alembert
equation $H(s+t)+H(s-t)=2H(s)H(t)$.  By the Acz\'el classification
theorem~\cite{Aczel1966}, the unique continuous solution with $H(0)=1$
and $H''(0)=1$ is $H(t)=\cosh t$, giving
$J(x)=\tfrac{1}{2}(x+x^{-1})-1$.
\end{proof}

\noindent The mirror symmetry $J(x) = J(x^{-1})$ (proved in
Theorem~\ref{thm:symmetry}) means neither chirality is
\emph{energetically} preferred.  The golden ratio $\varphi=(1+\sqrt{5})/2$,
forced by the RCL self-similarity, breaks this symmetry dynamically:
$\varphi^n \to \infty$ while $\varphi^{-n} \to 0$, providing an
exponential amplification cascade.

%% ============================================================
\section{Introduction}
\label{sec:intro}
%% ============================================================

Life universally uses L-amino acids in proteins and D-ribose in
DNA/RNA.  Why this specific chirality?  Recognition
Science~\cite{RS_Axioms} derives homochirality from the interaction
between $J$-cost symmetry and $\varphi$-cascade dynamics.  Three
competing hypotheses exist in the literature:
\begin{enumerate}
  \item \textbf{Random symmetry breaking}: Chance amplified by
  autocatalysis~\cite{Kondepudi1981}.
  \item \textbf{Circularly polarized light}: UV from neutron stars
  preferentially destroys D-amino acids.
  \item \textbf{Parity violation}: The weak nuclear force creates a
  tiny energy difference between enantiomers~\cite{Hegstrom1980}.
\end{enumerate}

None provides a zero-parameter first-principles derivation.

%% ============================================================
\section{$J$-Cost Symmetry}
\label{sec:symmetry}
%% ============================================================

\begin{theorem}[Cost Symmetry]
\label{thm:symmetry}
$J(x) = J(x^{-1})$ for all $x > 0$.
\end{theorem}

\begin{proof}
$J(x^{-1}) = \tfrac{1}{2}(x^{-1} + x) - 1
= \tfrac{1}{2}(x + x^{-1}) - 1 = J(x)$.
\end{proof}

\begin{corollary}[No Intrinsic Chirality Preference]
A configuration with ratio $x$ and its mirror image with ratio
$x^{-1}$ have identical $J$-cost.  There is no intrinsic cost
difference between L and D configurations.
\end{corollary}

%% ============================================================
\section{The $\varphi$-Symmetry Breaking}
\label{sec:breaking}
%% ============================================================

\subsection{The $\varphi$ Identity}

\begin{theorem}[Unique Chirality Operator]
\label{thm:chirality-op}
The golden ratio is the unique positive real number satisfying
\begin{equation}
  \varphi - \varphi^{-1} = 1.
\end{equation}
\end{theorem}

\begin{proof}
$\varphi = (1 + \sqrt{5})/2$ satisfies $\varphi^2 = \varphi + 1$,
hence $\varphi - 1 = 1/\varphi$, i.e.,
$\varphi - \varphi^{-1} = 1$.  To show uniqueness: if
$x - x^{-1} = 1$ for $x > 0$, then $x^2 - x - 1 = 0$, which has
unique positive root $x = (1 + \sqrt{5})/2 = \varphi$.
\end{proof}

\begin{remark}
The identity $\varphi - \varphi^{-1} = 1$ is the mathematical source
of chirality breaking.  Although $J(\varphi) = J(\varphi^{-1})$
(cost symmetry), the \emph{cascading dynamics} of $\varphi^n$ versus
$\varphi^{-n}$ are asymmetric.
\end{remark}

\subsection{Cascade Amplification}

\begin{theorem}[$\varphi$-Cascade Amplification]
\label{thm:cascade}
Starting from an equal mixture of configurations with ratios $\varphi$
and $\varphi^{-1}$, after $n$ recognition cycles:
\begin{align}
  x_n^L &= \varphi^n \to \infty, \\
  x_n^D &= \varphi^{-n} \to 0, \\
  \frac{x_n^L}{x_n^D} &= \varphi^{2n} \to \infty.
\end{align}
The L-configuration dominates exponentially.
\end{theorem}

\begin{proof}
$\varphi > 1$, so $\varphi^n$ is a strictly increasing sequence
diverging to $+\infty$.  $\varphi^{-n} = 1/\varphi^n \to 0$.
The ratio $\varphi^{2n}$ diverges.
\end{proof}

\begin{remark}[Identification of L/D with $\varphi$-branches]
The mathematical cascade is clear: one branch ($\varphi^n$) grows
without bound while the other ($\varphi^{-n}$) decays.  The
identification of the dominant $\varphi^n$ branch with L-amino acids
(rather than D) is not derived from RS first principles; it depends on
the sign of the initial bias (whether the weak-force correction $\epsilon$
is positive or negative in the RS ratio convention).  RS predicts that
\emph{one} chirality dominates universally; the identification with
L-amino acids specifically relies on the empirical fact that
parity-violating energy differences favor L-amino acids by
$\Delta E_{\rm weak} \approx 10^{-17}$~eV (measured;
see~\cite{Hegstrom1980}).  The cascade then amplifies this empirical
bias.
\end{remark}

%% ============================================================
\section{The Initial Bias}
\label{sec:bias}
%% ============================================================

\begin{theorem}[Weak Force Bias Amplification]
\label{thm:bias}
The weak nuclear force contributes a parity-violating energy
difference $\Delta E_{\mathrm{weak}} \approx 10^{-17}$~eV between
enantiomers.  In RS, this corresponds to a $J$-cost asymmetry:
\begin{equation}
  \delta J = J(\varphi + \epsilon) - J(\varphi)
  \approx J'(\varphi) \cdot \epsilon,
\end{equation}
where $\epsilon \sim \Delta E_{\mathrm{weak}} / E_{\mathrm{coh}}
\sim 10^{-16}$.  This tiny bias is amplified by the cascade
to a ratio $\varphi^{2n}$.
\end{theorem}

\begin{proof}
After $n$ polymerization steps, the ratio of L to D configurations is
$\varphi^{2n} \cdot (1 + \delta J / J(\varphi))^n$.  For $n \gg 1$,
the $\varphi^{2n}$ factor dominates regardless of the magnitude of
$\delta J$.
\end{proof}

\begin{remark}
Even without the weak force, any infinitesimal statistical
fluctuation would be amplified.  The weak force provides a
\emph{deterministic} bias that selects L-amino acids specifically.
\end{remark}

%% ============================================================
\section{The D-Sugar / L-Amino Acid Correlation}
\label{sec:correlation}
%% ============================================================

\begin{theorem}[Complementary Chiralities Forced]
\label{thm:complement}
Life uses L-amino acids and D-sugars because these are the
complementary pair with $\varphi \times \varphi^{-1} = 1$, minimizing
the total $J$-cost of protein--nucleic acid recognition.
\end{theorem}

\begin{proof}[Structural argument]
In the RS ledger, recognition between a protein and a nucleic acid
requires their effective ratios to satisfy $J(x_{\mathrm{protein}}) =
J(x_{\mathrm{nucleic}})$ (equal recognition cost, otherwise one partner
dominates and the interaction is unstable).  By $J$-symmetry
(Theorem~\ref{thm:symmetry}), this is satisfied when
$x_{\mathrm{nucleic}} = x_{\mathrm{protein}}^{-1}$.  The
$\varphi$-cascade (Section~\ref{sec:breaking}) drives proteins to
$x_{\mathrm{protein}} = \varphi^n$ (the dominant branch), so nucleic
acids must lie at $x_{\mathrm{nucleic}} = \varphi^{-n}$.  Since proteins
adopt the L-configuration ($\varphi^n$ branch) and nucleic acids adopt
the complementary branch, D-ribose is the uniquely forced sugar
configuration.
\end{proof}

\begin{remark}
The $J$-cost matching criterion $J(x) = J(x^{-1})$ is satisfied by
\emph{any} reciprocal pair, not just $(\varphi^n, \varphi^{-n})$.  The
specific selection of the $\varphi$-ladder is forced by the cascade
dynamics of Section~\ref{sec:breaking}: $\varphi$ is the unique
fixed-point attractor of the recognition updating map, so all stable
polymer configurations are $\varphi$-ladder states.  The L/D
correlation then follows from the requirement that the two polymers
occupy reciprocal rungs.  A fully rigorous derivation of this
  selection from the RS Hamiltonian would require solving the polymer
  folding problem within RS, which is an open problem.
\end{remark}

\begin{corollary}[Mixed Chirality Forbidden]
A cell with mixed L/D amino acids would have $J$-cost divergence in
protein folding, destabilizing all proteins simultaneously.
Homochirality is not merely preferred but \emph{required} for stable
molecular recognition.
\end{corollary}

%% ============================================================
\section{Experimental Evidence}
\label{sec:evidence}
%% ============================================================

\begin{center}
\renewcommand{\arraystretch}{1.3}
\begin{tabular}{@{}lp{9cm}@{}}
\toprule
\textbf{Observation} & \textbf{RS Interpretation} \\
\midrule
Murchison meteorite ($\sim\!8\%$ L-excess) &
Initial weak-force asymmetry before cascade amplification. \\
Circularly polarized UV ($\sim\!0.1\%$ excess) &
Secondary effect on top of the primary $\varphi$-cascade. \\
$^{14}$C $\beta$-decay (L-amino acid excess) &
The parity-violating $\epsilon$ that seeds the cascade. \\
\bottomrule
\end{tabular}
\end{center}

%% ============================================================
\section{Predictions}
\label{sec:predictions}
%% ============================================================

\begin{enumerate}
  \item \textbf{Universal homochirality}: All life in the universe
  (if it exists) will use one homochiral configuration.  Which one
  depends on initial conditions in each biosphere.

  \item \textbf{Racemization rate}: In isolation, L-amino acids
  slowly racemize because $J(x) = J(x^{-1})$.  Rate:
  $k_{\mathrm{racem}} \propto \varphi^{-n}$ where $n$ is the polymer
  chain length.

  \item \textbf{Mixed chirality non-viable}: A cell with mixed L/D
  amino acids would have $J$-cost divergence in protein folding.
  No viable organism can use both chiralities simultaneously.

  \item \textbf{Protein structure test}: In any homochiral polymer,
  the ratio of next-to-nearest to nearest amino acid interactions is
  $\varphi^2/\varphi = \varphi$, observable in protein structure
  statistics.
\end{enumerate}

%% ============================================================
\section{Conclusion}
\label{sec:conclusion}
%% ============================================================

Biological homochirality arises from two RS principles:
\begin{enumerate}
  \item \textbf{$J$-cost symmetry}: $J(\varphi) = J(\varphi^{-1})$
  ensures no intrinsic chirality preference.
  \item \textbf{$\varphi$-cascade asymmetry}:
  $\varphi^n \to \infty$ while $\varphi^{-n} \to 0$, so one
  chirality dominates through exponential amplification.
\end{enumerate}

The exact identity $\varphi - \varphi^{-1} = 1$ uniquely identifies
the golden ratio as the chirality-breaking operator.  The L/D
correlation (L-amino acids + D-sugars) is forced by complementary
recognition: $\varphi \times \varphi^{-1} = 1$ ensures proteins and
nucleic acids have equal $J$-cost and can bind each other.

%% ============================================================
\begin{thebibliography}{99}

\bibitem{Kondepudi1981}
D.~Kondepudi and I.~Prigogine, ``Sensitivity of Nonequilibrium
Chemical Systems to Parity Violation,'' Physica A \textbf{107}, 1
(1981).

\bibitem{Hegstrom1980}
R.~A.~Hegstrom, D.~Rein, and P.~G.~H.~Sandars, ``Calculation of
the Parity-Violating Energy Difference Between Mirror-Image
Molecules,'' J.\ Chem.\ Phys.\ \textbf{73}, 2329 (1980).

\bibitem{Blackmond2010}
D.~G.~Blackmond, ``The Origin of Biological Homochirality,''
Cold Spring Harbor Perspectives in Biology \textbf{2}, a002147
(2010).

\bibitem{RS_Axioms}
J.~Washburn, ``The Algebra of Reality: A Recognition Science Derivation
of Physical Law,'' \emph{Axioms} \textbf{15}(2), 90 (2026).
\texttt{doi:10.3390/axioms15020090}

\bibitem{Aczel1966}
J.~Acz\'el, \emph{Lectures on Functional Equations and Their
Applications} (Academic Press, 1966).

\end{thebibliography}

\end{document}
